\section{Discussion}
\label{sec:discussion}

\subsection{Implications of NP-hardness for the DIA}

Theorem~\ref{thm:nphard} establishes that no polynomial-time algorithm
can guarantee an optimal redistricting map unless P = NP.
P = NP is universally believed to be false among theoretical computer
scientists; making the NP-hardness result effectively unconditional for
all practical purposes.
This has two important implications for the Districting Integrity Act.

First, the DIA's use of a heuristic (METIS) is not a design weakness but
a theoretical necessity: any algorithm that claimed to find the optimal
compact map would be implicitly claiming to solve an NP-hard problem.
Second, legal challenges that argue the DIA should use the ``most compact''
possible algorithm are misdirected: finding the globally most compact map
is computationally intractable.

It is important to note that Theorem~\ref{thm:nphard} concerns the
\emph{optimisation} problem (finding the minimum-cut partition), not
feasibility (finding any valid partition).
A valid partition always exists for any connected graph with $k \leq n$:
any spanning tree can be cut into $k$ connected subtrees.
NP-hardness means that selecting the \emph{best} partition is
computationally intractable --- not that redistricting is impossible.

\subsection{Approximation quality in practice}

Proposition~\ref{prop:approx} establishes that no known polynomial-time
algorithm achieves better than $O(\log n)$ approximation for balanced
$k$-partition in general; the best known SDP-based algorithm
(ARV~\citep{arora2009expander}) achieves $O(\sqrt{\log n})$ for Sparsest
Cut.  No formal approximation ratio is claimed for METIS's heuristic.

Empirically, B.7~\citep{b7paper} shows that the DIA seed is within 3\%
of the best partition found across 10,000 independent seeds.
This empirical gap is not the same as the theoretical approximation ratio
relative to the true optimum (which is unknown); the reference point is
the best of 10,000 seeds, not a provably optimal solution.
The empirical 3\% gap should be interpreted as evidence of practical
solution quality, not as a formal bound.
Legal challenges should note that the worst-case guarantee from complexity
theory is much weaker than the empirical performance: METIS
systematically finds solutions close to the best feasible result found
by exhaustive seed search at the redistricting scale.

\subsection{Comparison to alternative algorithms}

Markov chain approaches~\citep{fifield2020} can sample from the solution
space and find near-optimal maps, but their convergence time is
exponential in the worst case.  IP-based approaches~\citep{mehrotra1998optimal}
can find exact solutions for small instances ($n \leq 500$) but do not
scale to census-tract graphs.

The METIS runtime of $O(n^{1.07})$ is the best available among algorithms
that produce legally compliant (contiguous, population-balanced) maps at
census-tract scale.  The next fastest alternatives (simulated annealing,
genetic algorithms) are $O(n^{1.5})$ or worse on the same problem class.

\subsection{Limitations}

The NP-hardness reduction in Theorem~\ref{thm:nphard} holds for the
optimisation version (find a partition with minimum edge cut subject to
connectivity and population balance).
The decision version (does a partition with edge cut $\leq C$ exist?)
is co-NP-hard by standard argument.
No formal approximation ratio for METIS recursive bisection is established
here; tighter bounds specific to the redistricting graph class (sparse
planar graphs with bounded diameter and integer population weights) remain
an open problem.

The empirical scaling ($b = 1.07$, $k > 1$ states) is fit on the 50-state
sample and may not generalise to block-level resolution
($n \approx 8.1\mathrm{M}$ nationally), which has not been tested at
full scale.
